\section*{Chapitre \ref{ch:b3:extension-strike-slip} --- Tectonique extensive et décrochante}

\begin{solution}{exo:b3:extension-strike-slip:1}
Les couches ont un pendage de $50 - 35 = \ang{15}$, à l'opposé des failles~; $\beta =
\sin\ang{50}/\sin\ang{35} = 0.766/0.574 = 1.34$.
\end{solution}

\begin{solution}{exo:b3:extension-strike-slip:2}
Relais dextre vers la droite~: bassin. Relais dextre vers la gauche~: crête.
Relais senestre vers la gauche~: bassin. Relais senestre vers la droite~: crête.
\end{solution}

\begin{solution}{exo:b3:extension-strike-slip:3}
$105/20 = \qty{5.25}{km/Ma} = \qty{5.25}{mm/an}$, proche de la
valeur de \qty{6}{mm/an}~: le taux n'a pas nécessairement été constant pendant \qty{20}{Ma}.
\end{solution}

\begin{solution}{exo:b3:extension-strike-slip:4}
\qty{200}{\percent} signifie $\Delta L = 2L_0$~: $L_0 = \qty{117.5}{km}$ et $L
= \qty{352.5}{km}$. Au nord, $\Delta L = 0.5L_0$~: $L_0 = \qty{470}{km}$
et $L = \qty{705}{km}$.
\end{solution}

\begin{solution}{exo:b3:extension-strike-slip:5}
$563/20 = 28$ à $563/15 = \qty{38}{mm/an}$. Le décalage est un minimum, et
la dérive mesurée par relevés géodésiques de part et d'autre de la zone de faille comprend aussi le mouvement sur
les failles voisines du système et la déformation élastique accumulée entre
les séismes.
\end{solution}

\begin{solution}{exo:b3:extension-strike-slip:6}
$t = 90 - \text{pendage}$. Pour \ang{60}~: $(1 + 0.6\times1.732)/(1 - 0.6\times
0.577) = 2.039/0.654 = 3.12$~; pour \ang{50}~: $1.715/0.497 = 3.45$~; pour
\ang{40}~: $1.504/0.285 = 5.28$. Blocage à $\arctan0.6 = \ang{31.0}$.
\end{solution}

\begin{solution}{exo:b3:extension-strike-slip:7}
$\arctan0.4 = \ang{21.8}$ et $\arctan0.2 = \ang{11.3}$. Une faille de pendage
\ang{15} ne glisse que si $\mu < \tan\ang{15} = 0.27$.
\end{solution}

\begin{solution}{exo:b3:extension-strike-slip:8}
$2.5\times4.0 = \qty{10.0}{km^2}$ et $2.5\times(4.0 - 1.5) =
\qty{6.25}{km^2}$.
\end{solution}

\begin{solution}{exo:b3:extension-strike-slip:9}
$132/1.19 = \qty{110.9}{km}$, étiré de \qty{21.1}{km}, ce qui demande
$21.1/6 = \qty{3.5}{Ma}$.
\end{solution}

\begin{solution}{exo:b3:extension-strike-slip:10}
$3.0/\sin\ang{60} = \qty{3.464}{km}$ avant et $3.0/\sin\ang{45} =
\qty{4.243}{km}$ après~: \qty{0.779}{km} par bloc ($\beta = 1.22$). Les couches
ont un pendage de $60 - 45 = \ang{15}$. Chaque bloc tourne d'un seul tenant~: la rotation qui
abaisse le pendage des failles plongeant, disons, vers l'est, fait tourner une ligne horizontale dans le
bloc de sorte qu'elle remonte vers l'est, et les couches plongent vers l'ouest, à l'opposé des
failles.
\end{solution}

\begin{solution}{exo:b3:extension-strike-slip:11}
Dans une zone senestre est--ouest, $\sigma_1$ s'oriente à \ang{045}. Les cisaillements R
se forment à $\phi/2 = \ang{20}$ de la zone, du côté de $\sigma_1$~: direction \ang{070},
senestre. R$'$ à $90 - 20 = \ang{70}$~: direction \ang{020}, dextre. Fractures de tension
parallèles à $\sigma_1$~: direction \ang{045}.
\end{solution}

\begin{solution}{exo:b3:extension-strike-slip:12}
Pour $\mu = 0.6$, les failles se forment à \ang{60.5} et se bloquent à \ang{31.0}~: $\beta =
\sin\ang{60.5}/\sin\ang{31.0} = 0.870/0.515 = 1.69$. Pour $\mu = 0.85$~: de
\ang{65.2} à \ang{40.4}, $\beta = 0.908/0.648 = 1.40$. Les failles à plus faible
frottement pivotent davantage avant de se bloquer~: les roches faibles s'étendent plus avant que
de nouvelles failles se forment.
\end{solution}

\begin{solution}{pb:b3:extension-strike-slip:1}
\textbf{1.} Les couches étaient horizontales et ont pivoté avec les blocs~; les
failles ont pivoté du même angle, donc $\theta_0 = 42 + 18 = \ang{60}$,
le pendage d'Anderson pour $\phi = \ang{30}$.

\textbf{2.} $\beta = \sin\ang{60}/\sin\ang{42} = 0.8660/0.6691 = 1.294$.

\textbf{3.} $24.0/1.294 = \qty{18.5}{km}$~; le domaine a été étendu de
\qty{5.5}{km} (\qty{5.46}{km}).

\textbf{4.} $t = \ang{48}$~: $(1 + 0.6\times0.900)/(1 - 0.6\times1.111) =
1.540/0.334 = 4.62$, contre $3.12$ pour de nouvelles failles à \ang{60.5}~: environ
moitié plus.

\textbf{5.} À $\arctan0.6 = \ang{31.0}$~; $\beta = \sin\ang{60}/\sin\ang{31}
= 1.68$ au plus. Avant cela, de nouvelles failles d'environ \ang{60} recouperaient les
blocs, et une seconde génération de \omterm{def:b3:extension-strike-slip:domino}{failles en dominos} prendrait le relais.

\textbf{6.} Chaque colonne du toit s'abaisse selon la chute verticale de la faille
située sous elle, rapportée au rejet horizontal~; la faille est la plus raide près de la surface et s'aplatit
en profondeur, de sorte que les colonnes proches de la faille s'abaissent le plus et les plus éloignées
le moins~: les couches basculent vers la faille.

\textbf{7.} $D = A/H = 16.0/3.2 = \qty{5.0}{km}$.

\textbf{8.} $D$ est proportionnel à $A$~: $\qty{5.0}{km} \pm \qty{0.5}{km}$.

\textbf{9.} $3.2\times(5.0 - 2.0) = \qty{9.6}{km^2}$.

\textbf{10.} Le rejet horizontal, \qty{3.2}{km}.

\textbf{11.} Le remplissage s'épaissit vers la faille, en biseau, et ses couches les plus
anciennes sont plus basculées que les plus jeunes~: chaque couche s'est déposée pendant que
le roll-over et l'espace voisin de la faille continuaient de s'agrandir.

\textbf{12.} $z = P/(\rho g) = \num{5.5e8}/(2800\times9.8) = \qty{20.0}{km}$.

\textbf{13.} $20.0/(24 - 16) = \qty{2.5}{km/Ma} = \qty{2.5}{mm/an}$.

\textbf{14.} $20.0/\sin\ang{30} = \qty{40}{km}$ de glissement, et $40\cos\ang{30} =
\qty{34.6}{km}$ d'extension horizontale.

\textbf{15.} Non~: \ang{15} est inférieur au pendage de verrouillage de \ang{31}. Soit la
faille a glissé pendant qu'elle était plus pentue, puis a été basculée (\omterm{def:b3:extension-strike-slip:core-complex}{charnière roulante}), soit
elle était faible, avec $\mu < 0.27$ dus à des argiles ou à des fluides à haute pression.

\textbf{16.} Par $45 - 15 = \ang{30}$~: lorsque la plaque supérieure s'est éloignée, le
mur déchargé s'est soulevé par isostasie et s'est courbé, faisant pivoter la faille
à travers la charnière.

\textbf{17.} \qty{34.6}{km}, contre \qty{8.7}{km} pour le demi-graben (question
23)~: quatre fois plus. Les core complexes se situent là où l'extension est la plus forte, comme dans
le Basin and Range central.

\textbf{18.} Une faille senestre en relais vers la gauche~: \omterm{def:b3:extension-strike-slip:bends}{coude extensif}, un bassin s'ouvre.

\textbf{19.} $\qty{5}{mm/an}\times\qty{1.2}{Ma} = \qty{6.0}{km}$.

\textbf{20.} $4.0\times6.0 = \qty{24}{km^2}$~; $3.0 + 6.0 = \qty{9.0}{km}$ de longueur~;
$9.0/4.0 = 2.25$.

\textbf{21.} $132/10 = 13$ à $132/7 = 19$~: le bassin de la mer Morte est bien plus
allongé, son décalage vaut plusieurs fois sa largeur.

\textbf{22.} Un \omterm{def:b3:extension-strike-slip:bends}{coude compressif}~: une crête soulevée entre les branches~; en coupe,
une \omterm{def:b3:extension-strike-slip:flower}{structure en fleur} positive, la faille se divisant vers le haut en branches
à composante inverse, comme celle de 100 à \qty{300}{m} de large sur la
faille Dead Sea Transform.

\textbf{23.} $3.2 + 5.46 = \qty{8.66}{km}$~: le demi-graben de Merrow a subi
une extension de \textbf{\qty[text-series-to-math=true]{8.7}{km}}.
\end{solution}
